\chapter{Conclusion}
\label{ch:conclusion}

This dissertation has examined the short-run pass-through of global oil price shocks to India's headline CPI inflation over April 2004 to December 2024, with a focus on whether the pass-through is asymmetric, that is, whether oil price increases raise inflation by more than decreases lower it.

The main findings can be summarised as follows. First, positive oil price shocks are associated with higher monthly CPI inflation. The cumulative positive pass-through estimate of 0.021 implies that a 10 per cent increase in rupee-denominated oil prices raises monthly inflation by about 0.21 percentage points. This is an economically meaningful effect, even if individual lag coefficients are not all statistically significant.

Second, the estimated pass-through from negative oil shocks is near zero, meaning that oil price declines do not appear to provide measurable relief to headline CPI inflation in the short run. The asymmetry gap, the difference between positive and negative cumulative pass-through, is positive across all specifications examined.

Third, however, the formal Wald test for the equality of positive and negative pass-through does not reject symmetry at the 5 per cent level ($p = 0.241$). The point estimates suggest asymmetry, but the statistical evidence is not strong enough to confirm it in headline CPI. This is consistent with the high noise-to-signal ratio in India's CPI basket, where food and non-oil components dominate.

Fourth, separating the Brent oil price from the rupee exchange rate improves the model fit and produces a marginally significant positive oil pass-through ($p = 0.093$). The exchange rate is independently significant ($p = 0.029$), confirming that rupee depreciation is a distinct inflationary channel in India.

Fifth, the Fuel and Light CPI sub-index shows a much stronger positive oil pass-through (0.061, $p = 0.034$), consistent with the expectation that a more energy-exposed price index should exhibit clearer oil transmission.

\section{Limitations}

Several limitations should be noted. The study uses headline CPI, which includes many components unrelated to oil, making it difficult to isolate the oil-specific signal. The IIP control variable is an imperfect proxy for economic activity. The analysis covers only the short run (up to three months of lags) and does not estimate long-run cointegrating relationships. The Fuel and Light appendix uses a shorter sample starting from 2011, which limits comparison with the main model.

\section{Directions for Further Research}

Extending the analysis to longer lag horizons using local projections or impulse response methods would help clarify whether the asymmetry becomes more pronounced over time. Disaggregated analysis across individual CPI sub-groups (transport, food, manufactured goods) could identify the specific channels through which oil prices affect headline inflation. An analysis incorporating inflation expectations data from the RBI's household survey would address whether oil price shocks affect expected inflation differently when they are increases versus decreases. Finally, as more post-deregulation data become available, revisiting the sub-sample comparison with a longer post-2014 window may yield sharper estimates of the deregulation effect on pass-through symmetry.
